% GENERATED FILE. Edit canonical lessons, then run npm run generate:books.
% canonical-source-hash: fnv1a64:11ca9894d53fd138
% canonical-lessons: LA-C320-pavor, LA-C320-metus, LA-C320-formido, LA-C320-pavidus, LA-C320-trepido

\chapter{Pavor}
\label{ch:la-pavor}
\hlchaptermodality{320}

\begin{chapteropening}
\textbf{By the end of this chapter:} \emph{I can describe fear in a story: panic, fear, dread, frightened and to be in a flurry.}

Complete the last lesson of chapter 320: I can describe fear in a story: panic, fear, dread, frightened and to be in a flurry.
\end{chapteropening}

\section[pavor, pavōris]{pavor, pavōris — panic, terror}

\label{lesson:LA-C320-pavor}



\begin{quote}
Before the new one: say the Latin for the watch, then the Latin for to scorch.
\end{quote}

\subsection*{You'll want to know: pavor, pavōris}
\textbf{pavor, pavōris} — \textquotedblleft{}panic, terror\textquotedblright{}.

\begin{quote}
Pavor omnēs cīvēs cēpit.
\end{quote}

\begin{quote}
\emph{Panic seized all the citizens.}
\end{quote}

Masculine, 3rd declension. A sudden fear that shakes the body.

\begin{cousinweb}
\textbf{pavor} is built on \emph{paveō}, \textquotedblleft{}to tremble with fear\textquotedblright{} (EtymDB): the trembling itself.
\end{cousinweb}

\begin{grammarlens}[title={What you've built}]
One more word for telling, explaining and describing at B1.
\end{grammarlens}

\subsection*{Your turn}
\begin{itemize}
\raggedright
  \item \emph{Say it:} \emph{pavor, pavōris}
  \item \emph{Say it:} \emph{pavor, pavōris}, once more
  \item \emph{Say it:} a sentence of your own with \emph{pavor, pavōris}
  \item \emph{Recall:} say the Latin for the watch, then the Latin for to scorch, then say \emph{pavor, pavōris} again
\end{itemize}

\begin{tcolorbox}[breakable,colback=teal!4,colframe=teal!35!black,title={Before you move on}]
What does pavor, pavōris mean? (\textquotedblleft{}Panic, terror\textquotedblright{}.) Say it once more.
\end{tcolorbox}
\section[metus, metūs]{metus, metūs — fear, dread}

\label{lesson:LA-C320-metus}



\begin{quote}
Before the new one: say the Latin for to scorch, then the Latin for panic.
\end{quote}

\subsection*{You'll want to know: metus, metūs}
\textbf{metus, metūs} — \textquotedblleft{}fear, dread\textquotedblright{}.

\begin{quote}
Propter metum hostium portās clausērunt.
\end{quote}

\begin{quote}
\emph{Out of fear of the enemy they shut the gates.}
\end{quote}

Masculine, 4th declension. A steady, lasting fear, beside \emph{timor} (taught).

\begin{cousinweb}
\textbf{metus} is a Latin word whose earlier origin is not settled (Lewis \& Short, EtymDB). Its verb is \emph{metuō}, \textquotedblleft{}to fear\textquotedblright{} (taught).
\end{cousinweb}

\begin{grammarlens}[title={What you've built}]
One more word for telling, explaining and describing at B1.
\end{grammarlens}

\subsection*{Your turn}
\begin{itemize}
\raggedright
  \item \emph{Say it:} \emph{metus, metūs}
  \item \emph{Say it:} \emph{metus, metūs}, once more
  \item \emph{Say it:} a sentence of your own with \emph{metus, metūs}
  \item \emph{Recall:} say the Latin for to scorch, then the Latin for panic, then say \emph{metus, metūs} again
\end{itemize}

\begin{tcolorbox}[breakable,colback=teal!4,colframe=teal!35!black,title={Before you move on}]
What does metus, metūs mean? (\textquotedblleft{}Fear, dread\textquotedblright{}.) Say it once more.
\end{tcolorbox}
\section[formīdō, formīdinis]{formīdō, formīdinis — dread, terror}

\label{lesson:LA-C320-formido}



\begin{quote}
Before the new one: say the Latin for panic, then the Latin for fear.
\end{quote}

\subsection*{You'll want to know: formīdō, formīdinis}
\textbf{formīdō, formīdinis} — \textquotedblleft{}dread, terror\textquotedblright{}.

\begin{quote}
Magna formīdō mīlitēs tenēbat.
\end{quote}

\begin{quote}
\emph{Great dread held the soldiers.}
\end{quote}

Feminine, 3rd declension. There is also a verb \emph{formīdō, formīdāre}, 'to dread'; English \emph{formidable} comes from it.

\begin{cousinweb}
\textbf{formīdō} is explained by Lewis \& Short as \textquotedblleft{}the fear that makes one rigid\textquotedblright{}, from a root also seen in \emph{firmus}, \textquotedblleft{}firm\textquotedblright{}; the link is not certain. English \emph{formidable} comes from its family.
\end{cousinweb}

\begin{grammarlens}[title={What you've built}]
One more word for telling, explaining and describing at B1.
\end{grammarlens}

\subsection*{Your turn}
\begin{itemize}
\raggedright
  \item \emph{Say it:} \emph{formīdō, formīdinis}
  \item \emph{Say it:} \emph{formīdō, formīdinis}, once more
  \item \emph{Say it:} a sentence of your own with \emph{formīdō, formīdinis}
  \item \emph{Recall:} say the Latin for panic, then the Latin for fear, then say \emph{formīdō, formīdinis} again
\end{itemize}

\begin{tcolorbox}[breakable,colback=teal!4,colframe=teal!35!black,title={Before you move on}]
What does formīdō, formīdinis mean? (\textquotedblleft{}Dread, terror\textquotedblright{}.) Say it once more.
\end{tcolorbox}
\section[pavidus, pavida, pavidum]{pavidus, pavida, pavidum — frightened, trembling}

\label{lesson:LA-C320-pavidus}



\begin{quote}
Before the new one: say the Latin for fear, then the Latin for dread.
\end{quote}

\subsection*{You'll want to know: pavidus, pavida, pavidum}
\textbf{pavidus, pavida, pavidum} — \textquotedblleft{}frightened, trembling\textquotedblright{}.

\begin{quote}
Pavida ancilla in angulō sedēbat.
\end{quote}

\begin{quote}
\emph{The frightened maid was sitting in the corner.}
\end{quote}

From the same verb as \emph{pavor}. Plautus pairs it with \emph{timidus} (taught).

\begin{cousinweb}
\textbf{pavidus} is built on \emph{paveō}, \textquotedblleft{}to tremble with fear\textquotedblright{}, with the adjective ending \emph{-idus} (EtymDB).
\end{cousinweb}

\begin{grammarlens}[title={What you've built}]
One more word for telling, explaining and describing at B1.
\end{grammarlens}

\subsection*{Your turn}
\begin{itemize}
\raggedright
  \item \emph{Say it:} \emph{pavidus, pavida, pavidum}
  \item \emph{Say it:} \emph{pavidus, pavida, pavidum}, once more
  \item \emph{Say it:} a sentence of your own with \emph{pavidus, pavida, pavidum}
  \item \emph{Recall:} say the Latin for fear, then the Latin for dread, then say \emph{pavidus, pavida, pavidum} again
\end{itemize}

\begin{tcolorbox}[breakable,colback=teal!4,colframe=teal!35!black,title={Before you move on}]
What does pavidus, pavida, pavidum mean? (\textquotedblleft{}Frightened, trembling\textquotedblright{}.) Say it once more.
\end{tcolorbox}
\section[trepidō, trepidāre]{trepidō, trepidāre — to panic, to be in a flurry}

\label{lesson:LA-C320-trepido}



\begin{quote}
Before the new one: say the Latin for dread, then the Latin for frightened.
\end{quote}

\subsection*{You'll want to know: trepidō, trepidāre}
\textbf{trepidō, trepidāre} — \textquotedblleft{}to panic, to be in a flurry\textquotedblright{}.

\begin{quote}
Cīvēs in viīs trepidābant, nec quid facerent sciēbant.
\end{quote}

\begin{quote}
\emph{The citizens were rushing about in alarm in the streets and did not know what to do.}
\end{quote}

\emph{trepidō, trepidāre, trepidāvī, trepidātum}: anxious hurrying, the bustle of fear. English \emph{trepidation}.

\begin{cousinweb}
\textbf{trepidō} is built on \emph{trepidus}, \textquotedblleft{}restless, alarmed\textquotedblright{} (EtymDB). English \emph{trepidation} comes from it.
\end{cousinweb}

\begin{grammarlens}[title={What you've built}]
One more word for telling, explaining and describing at B1.
\end{grammarlens}

\subsection*{Your turn}
\begin{itemize}
\raggedright
  \item \emph{Say it:} \emph{trepidō, trepidāre}
  \item \emph{Say it:} \emph{trepidō, trepidāre}, once more
  \item \emph{Say it:} a sentence of your own with \emph{trepidō, trepidāre}
  \item \emph{Recall:} say the Latin for dread, then the Latin for frightened, then say \emph{trepidō, trepidāre} again
\end{itemize}

\begin{tcolorbox}[breakable,colback=teal!4,colframe=teal!35!black,title={Before you move on}]
What does trepidō, trepidāre mean? (\textquotedblleft{}To panic, to be in a flurry\textquotedblright{}.) Say it once more.
\end{tcolorbox}
